\section{Reproducing this document (toolchain)}\label{sec:repro}

This whitepaper is itself a small, self-contained build: a contributor who clones the
repository should be able to regenerate the PDF, with every diagram, from the sources
checked into \srcf{docs/ArchitectureReviewLatex/}. The layout and style contract are
inherited from the Part Placement design whitepaper
(\srcf{docs/PartPlacementDesignLatex/}), so the two documents read as one series; this
appendix records the small ways this build differs and how the review itself was produced.

\subsection{Directory layout and pipeline}

The master file \srcf{ArchitectureReview.tex} carries only the document scaffolding ---
title block, front matter, and an ordered list of \cppinline|\input| directives --- while
the prose lives in \srcf{sections/}, one file per section. The shared style contract is
isolated in \srcf{preamble.tex}: the color palette, the fonts (XCharter text,
\code{newtxmath}, Inconsolata mono), the callout environments, and the domain macros
(\code{\textbackslash src}, \code{\textbackslash srcf}, \code{\textbackslash tikzfig}, the
feature-status markers) that the section files depend on. A section file is just a body,
never a standalone \LaTeX{} program.

Unlike the Part Placement whitepaper, this document has no PlantUML stage: every figure is
hand-built Ti\emph{k}Z under \srcf{diagrams/tikz/}, included directly by
\code{\textbackslash tikzfig} and compiled in line by the main engine. The Makefile is
correspondingly simpler --- no Java, no \code{rsvg-convert}, no fetched jar:

\begin{minted}{text}
make            # latexmk -pdf -shell-escape ArchitectureReview.tex
make clean      # remove LaTeX aux files and the minted cache
make veryclean  # clean, and also remove ArchitectureReview.pdf
\end{minted}

The document engine is \code{pdflatex}, driven by \code{latexmk} so the multi-pass dance
(cross-references, the table of contents, \code{cleveref}, the bookmark tree) converges
automatically. The one non-default requirement is \code{-shell-escape}: listings are
typeset by \code{minted}, which shells out to \code{pygmentize} for syntax highlighting.
Prerequisites are a \TeX{} distribution providing \code{latexmk}, \code{pdflatex}, and
\code{minted} with its \code{pygmentize} helper.

\subsection{How this review was produced}

The review is pinned to one tree state: commit \code{20eca72}, the \code{development} tip
on 2026-07-15. It was originally produced against commit \code{254cd41} (``Major comment
cleanup'') on 2026-07-06 and revised in place after the \code{core/} restructure, the
\code{dynamic-Cd} merge, and the \code{RocketTreeViewImplementation} merge landed. Every
\code{\textbackslash src} citation refers to the pinned revision, with source paths written
include-rooted (\srcf{model/}, \srcf{sim/}, and \srcf{utils/} live under \srcf{core/} on
disk); line numbers will drift as the tree moves, so a reader on a later commit should
treat them as anchors, not absolutes.

The analysis behind the document was performed as a fan-out of independent subsystem
reviews --- part tree and placement, motors, simulation core, aerodynamics, GUI and
visualizer, CLI and persistence, build/test infrastructure, and a dedicated
incompleteness sweep --- each grounded in a full read of its files, followed by an
independent adversarial verification pass that re-checked every citation and status claim
against the source. The OpenRocket comparison baseline (\cref{sec:comparison}) was
researched from primary sources (\texttt{openrocket.info}, the OpenRocket documentation
and technical papers, and the \texttt{openrocket/openrocket} repository), then reconciled
against the verified QtRocket findings. Documents under \srcf{docs/} were deliberately
treated as claims to test, not facts to inherit: where a design document and the code
disagree, this review reports the code and flags the discrepancy.

Two behaviors were additionally confirmed by execution rather than reading: the CLI
command inventory of \cref{sec:interfaces} was captured from a live
\code{build/cli/qtrocket-cli} session, and end-to-end flight behavior was exercised through
the same REPL against a bundled design.
